\documentclass[12pt,a4paper]{article}
\usepackage[utf8]{inputenc}
\usepackage[T1]{fontenc}
\usepackage{amsmath,amssymb}
\usepackage{graphicx}
\usepackage{booktabs}
\usepackage{array}
\usepackage{caption}
\usepackage{hyperref}
\usepackage{geometry}
\usepackage{fancyhdr}
\usepackage{setspace}
\usepackage{xcolor}
\usepackage{float}
\usepackage{framed}

\setlength{\headheight}{14pt}
\geometry{a4paper,top=2.5cm,bottom=2.5cm,left=2.5cm,right=2.5cm}
\pagestyle{fancy}
\fancyhf{}
\rhead{\small Análisis numérico de EDPs · Proyecto Final}
\lhead{\small Pista B: Problema Inverso con PINNs}
\cfoot{\thepage}
\renewcommand{\headrulewidth}{0.4pt}
\hypersetup{colorlinks=true,linkcolor=blue!70!black,citecolor=blue!70!black}
\onehalfspacing

\newcommand{\lhat}{\hat{\lambda}}
\newcommand{\ltrue}{\lambda_{\text{real}}}
\newcommand{\Lf}{\mathcal{L}_f}
\newcommand{\Lb}{\mathcal{L}_b}
\newcommand{\Li}{\mathcal{L}_{ic}}
\newcommand{\Ld}{\mathcal{L}_d}

% Comando para espacio de figura con etiqueta
\newcommand{\figspace}[2]{%
  \begin{figure}[H]
    \centering
    \fbox{\begin{minipage}{0.92\textwidth}
      \centering\vspace{#1}
      \textbf{\large INSERTAR: #2}
      \vspace{#1}
    \end{minipage}}
    \caption{CAPTION DE: #2}
    \label{fig:#2}
  \end{figure}
}

\begin{document}

\begin{titlepage}
  \centering\vspace*{2cm}
  {\Large\textbf{Análisis Numérico de Ecuaciones Diferenciales}\par}
  \vspace{0.5cm}{\large Proyecto Final\par}\vspace{2cm}
  {\LARGE\textbf{Redes Neuronales Informadas por la Física (PINNs):}\par}
  \vspace{0.4cm}{\LARGE\textbf{Problema Inverso de Parámetros}\par}
  \vspace{0.5cm}{\large Pista B\par}\vspace{3cm}
  {\large\textit{Ecuación de difusión 1D con coeficiente desconocido}\par}
  \vspace{3cm}
  {\normalsize Basado en: Raissi, Perdikaris y Karniadakis (2019),\\
  \textit{J. Comput. Phys.} 378:686--707\par}
  \vfill{\normalsize\today\par}
\end{titlepage}

\tableofcontents\newpage

\section{Introducción y Planteamiento del Problema}

En numerosos fenómenos de transporte, la forma funcional de la EDP es conocida, pero
ciertos coeficientes constitutivos permanecen inaccesibles a medición directa. El
\textbf{problema inverso de parámetros} consiste en inferirlos a partir de observaciones
dispersas y ruidosas del campo solución.

Se estudia la \textbf{ecuación de difusión unidimensional}:
\begin{equation}
  \frac{\partial u}{\partial t}=\lambda\,\frac{\partial^2 u}{\partial x^2},\qquad
  (x,t)\in[0,1]^2,\label{eq:difusion}
\end{equation}
con condición inicial $u(x,0)=\sin(\pi x)$ y condiciones Dirichlet homogéneas
$u(0,t)=u(1,t)=0$. La solución analítica exacta es
\begin{equation}
  u(x,t)=\sin(\pi x)\,e^{-\lambda\pi^2 t}.\label{eq:analitica}
\end{equation}
El parámetro $\lambda>0$ es \textbf{incógnita}; el objetivo es recuperarlo a partir de
$N_d=100$ observaciones con 1\,\% de ruido gaussiano.

%──────────────────────────────────────────────────────────────────────────────
% FIGURA 1
%──────────────────────────────────────────────────────────────────────────────
\begin{figure}[H]
  \centering
  \fbox{\begin{minipage}{0.6\textwidth}\centering\vspace{2.5cm}
    \textbf{FIGURA 1 — Condición inicial $\sin(\pi x)$}
  \vspace{2.5cm}\end{minipage}}
  \caption{Condición inicial $u(x,0)=\sin(\pi x)$ utilizada en todos los experimentos.}
  \label{fig:ic}
\end{figure}

%──────────────────────────────────────────────────────────────────────────────
% FIGURA 2
%──────────────────────────────────────────────────────────────────────────────
\begin{figure}[H]
  \centering
  \fbox{\begin{minipage}{0.6\textwidth}\centering\vspace{2.5cm}
    \textbf{FIGURA 2 — Datos observados (100 puntos con 1\% de ruido)}
  \vspace{2.5cm}\end{minipage}}
  \caption{Distribución de los $N_d=100$ puntos de observación en el dominio
           $(x,t)\in[0,1]^2$, coloreados por el valor de $u$ observado con 1\,\% de ruido gaussiano.}
  \label{fig:datos}
\end{figure}

\section{Metodología}

\subsection{Arquitectura de la red}
Se implementó una PINN como subclase de \texttt{tf.keras.Model} en TensorFlow~2.x:
3~capas ocultas de 40~neuronas, activación $\tanh$, inicialización Glorot normal. El
coeficiente $\lhat$ se declara como \texttt{tf.Variable} escalar con restricción
$\lhat\geq10^{-6}$, inicializado en $\lhat_0=0.5$.

\subsection{Función de pérdida}
\begin{equation}
  \mathcal{L}(\theta,\lhat)=w_f\Lf+\Li+\Lb+w_d\Ld,\quad w_f=1,\;w_d=10.
\end{equation}

\subsection{Diferenciación automática}
Se usaron tapes anidados (\texttt{tape1} interno y \texttt{tape2} externo) para calcular
correctamente $\partial^2\tilde{u}/\partial x^2$.

\subsection{Entrenamiento}
Adam, $\eta=10^{-3}$, 8000~épocas. $\lhat$ y $\theta$ se optimizan conjuntamente.

\section{Problema Directo}

\subsection{Convergencia}

%──────────────────────────────────────────────────────────────────────────────
% FIGURA 3
%──────────────────────────────────────────────────────────────────────────────
\begin{figure}[H]
  \centering
  \fbox{\begin{minipage}{0.9\textwidth}\centering\vspace{2cm}
    \textbf{FIGURA 3 — Convergencia problema directo (pérdida total, PDE, IC+BC)}
  \vspace{2cm}\end{minipage}}
  \caption{Convergencia del problema directo: pérdida total, residuo PDE y
           condiciones de contorno/iniciales en escala logarítmica.}
  \label{fig:conv_directo}
\end{figure}

\subsection{Solución reconstruida}

%──────────────────────────────────────────────────────────────────────────────
% FIGURA 4
%──────────────────────────────────────────────────────────────────────────────
\begin{figure}[H]
  \centering
  \fbox{\begin{minipage}{0.95\textwidth}\centering\vspace{2cm}
    \textbf{FIGURA 4 — Solución directa: Analítica | PINN | Error absoluto (mapa de color)}
  \vspace{2cm}\end{minipage}}
  \caption{Problema directo: solución analítica (izquierda), solución PINN (centro) y
           error absoluto (derecha) en el dominio $[0,1]^2$.}
  \label{fig:sol_directo}
\end{figure}

\section{Problema Inverso}

\subsection{Convergencia del parámetro}

%──────────────────────────────────────────────────────────────────────────────
% FIGURA 5
%──────────────────────────────────────────────────────────────────────────────
\begin{figure}[H]
  \centering
  \fbox{\begin{minipage}{0.85\textwidth}\centering\vspace{2cm}
    \textbf{FIGURA 5 — Evolución conjunta de $\hat\lambda$ y pérdida total (doble eje)}
  \vspace{2cm}\end{minipage}}
  \caption{Evolución conjunta de $\lhat$ y la pérdida total durante el entrenamiento del
           problema inverso. La línea roja discontinua indica $\ltrue=0.1$.}
  \label{fig:conv_lambda}
\end{figure}

%──────────────────────────────────────────────────────────────────────────────
% FIGURA 6
%──────────────────────────────────────────────────────────────────────────────
\begin{figure}[H]
  \centering
  \fbox{\begin{minipage}{0.85\textwidth}\centering\vspace{2cm}
    \textbf{FIGURA 6 — Pérdida total problema inverso (escala log)}
  \vspace{2cm}\end{minipage}}
  \caption{Evolución de la pérdida total durante el entrenamiento del problema inverso
           en escala logarítmica.}
  \label{fig:loss_inv}
\end{figure}

\subsection{Resultado final}

\begin{table}[H]
  \centering
  \caption{Resultado del problema inverso.}\label{tab:resultado}
  \begin{tabular}{lc}\toprule
    Magnitud & Valor\\\midrule
    $\ltrue$ & 0.10000\\
    $\lhat$ estimado & 0.10017\\
    Error relativo & \textbf{0.17\,\%}\\
    $|\nabla_{\lhat}\mathcal{L}|$ final & $1.58\times10^{-6}$\\\bottomrule
  \end{tabular}
\end{table}

\subsection{Reconstrucción del campo}

%──────────────────────────────────────────────────────────────────────────────
% FIGURA 7
%──────────────────────────────────────────────────────────────────────────────
\begin{figure}[H]
  \centering
  \fbox{\begin{minipage}{0.95\textwidth}\centering\vspace{2cm}
    \textbf{FIGURA 7 — Campo inverso: PINN ($\hat\lambda=0.10017$) | Analítica | Error absoluto}
  \vspace{2cm}\end{minipage}}
  \caption{Problema inverso: solución PINN con $\lhat=0.10017$ (izquierda), solución
           analítica (centro) y error absoluto (derecha).}
  \label{fig:sol_inverso}
\end{figure}

\section{Validación}

\subsection{Comparación con diferencias finitas (Crank-Nicolson)}

\begin{table}[H]
  \centering
  \caption{Errores $L^2$ relativos de la validación contra Crank-Nicolson.}\label{tab:cn}
  \begin{tabular}{lcc}\toprule
    Comparación & Problema Directo & Problema Inverso ($\lhat=0.10017$)\\\midrule
    CN vs.\ analítica   & $6.711\times10^{-3}$ & $6.737\times10^{-3}$\\
    PINN vs.\ analítica & $1.470\times10^{-3}$ & $1.565\times10^{-3}$\\
    CN vs.\ PINN        & $7.453\times10^{-3}$ & $7.250\times10^{-3}$\\\bottomrule
  \end{tabular}
\end{table}

%──────────────────────────────────────────────────────────────────────────────
% FIGURA 8
%──────────────────────────────────────────────────────────────────────────────
\begin{figure}[H]
  \centering
  \fbox{\begin{minipage}{0.95\textwidth}\centering\vspace{1.8cm}
    \textbf{FIGURA 8 — Cortes temporales CN: problema DIRECTO\\
    (5 paneles: $t=0,\,0.25,\,0.5,\,0.75,\,1.0$; curvas Analítica, CN, PINN)}
  \vspace{1.8cm}\end{minipage}}
  \caption{Problema directo: cortes temporales en $t\in\{0,0.25,0.5,0.75,1.0\}$
           para la solución analítica, Crank-Nicolson y PINN.}
  \label{fig:cn_directo}
\end{figure}

%──────────────────────────────────────────────────────────────────────────────
% FIGURA 9
%──────────────────────────────────────────────────────────────────────────────
\begin{figure}[H]
  \centering
  \fbox{\begin{minipage}{0.95\textwidth}\centering\vspace{1.8cm}
    \textbf{FIGURA 9 — Cortes temporales CN: problema INVERSO\\
    (5 paneles: $t=0,\,0.25,\,0.5,\,0.75,\,1.0$; curvas Analítica, CN, PINN Inverso;\\
    título incluye $\hat\lambda=0.10017$)}
  \vspace{1.8cm}\end{minipage}}
  \caption{Problema inverso: cortes temporales en $t\in\{0,0.25,0.5,0.75,1.0\}$
           para la solución analítica, Crank-Nicolson y PINN ($\lhat=0.10017$).}
  \label{fig:cn_inverso}
\end{figure}

\subsection{Estudio de sensibilidad}

\begin{table}[H]
  \centering
  \caption{Error relativo en $\lhat$ (\%) — media sobre tres semillas (rango entre paréntesis).}
  \label{tab:sensibilidad}
  \begin{tabular}{ccc}\toprule
    $N_d$ & Ruido 0\,\% & Ruido 1\,\%\\\midrule
    20  & 0.18\,\% (0.03--0.26) & 0.35\,\% (0.11--0.52)\\
    50  & 0.23\,\% (0.17--0.29) & 0.12\,\% (0.07--0.21)\\
    100 & 0.21\,\% (0.10--0.29) & 0.65\,\% (0.53--0.80)\\
    200 & 0.19\,\% (0.14--0.28) & 0.26\,\% (0.07--0.65)\\\bottomrule
  \end{tabular}
\end{table}

%──────────────────────────────────────────────────────────────────────────────
% FIGURA 10
%──────────────────────────────────────────────────────────────────────────────
\begin{figure}[H]
  \centering
  \fbox{\begin{minipage}{0.95\textwidth}\centering\vspace{2cm}
    \textbf{FIGURA 10 — Estudio de sensibilidad\\
    (2 paneles: error en $\hat\lambda$ vs $N_d$ | error $L^2$ vs $N_d$;\\
    dos curvas por panel: Ruido 0\% y Ruido 1\%; línea umbral 1\%)}
  \vspace{2cm}\end{minipage}}
  \caption{Estudio de sensibilidad: error relativo en $\lhat$ (izquierda) y error $L^2$
           relativo del campo (derecha) en función de $N_d$, para
           $\sigma\in\{0\%,1\%\}$. Línea gris: umbral del 1\,\%.}
  \label{fig:sens}
\end{figure}

\subsection{Validación de plausibilidad física}
\begin{enumerate}
  \item \textbf{Positividad:} $\lhat\leftarrow\max(\lhat,10^{-6})$ tras cada paso.
  \item \textbf{Consistencia del campo:} reproduce el decaimiento exponencial analítico.
  \item \textbf{Diagnóstico de gradiente:} $|\nabla_{\lhat}\mathcal{L}|$ consistente durante todo el entrenamiento.
\end{enumerate}

\section{Discusión y Limitaciones}

\subsubsection*{Paisaje de optimización no convexo.}
Se detectó un mínimo local espurio ($N_d=50$, ruido 0\,\%, error $L^2=4.12\times10^{-2}$)
no reproducido en corridas posteriores, lo que motiva el uso de múltiples inicializaciones.

\subsubsection*{Ausencia de tendencia monótona con $N_d$.}
El error en $\lhat$ oscila entre 0.18\,\% y 0.23\,\% para ruido nulo, indicando que el
cuello de botella es el optimizador, no la cantidad de datos.

\subsubsection*{Variabilidad con ruido.}
Con ruido del 1\,\%, la variabilidad inter-semilla aumenta, especialmente en $N_d=100$.

\subsubsection*{Extensión a Fokker-Planck.}
Extensión natural: inferir deriva y difusión conjuntamente siguiendo Chen et al.\ (2021).

\section{Conclusiones}
\begin{enumerate}
  \item $\lambda$ recuperado con error relativo \textbf{0.17\,\%} (umbral: 1\,\%).
  \item Error $L^2$ del campo: $1.57\times10^{-3}$ (objetivo: $\sim10^{-3}$).
  \item PINN supera a Crank-Nicolson ($N_x=99$, $N_t=1000$) en precisión.
  \item Robusto para $N_d\geq20$ y ruido hasta 1\,\% — 24/24 configuraciones bajo umbral.
\end{enumerate}

\begin{thebibliography}{9}
\bibitem{raissi2019} Raissi, M., Perdikaris, P., \& Karniadakis, G.\,E. (2019).
Physics-informed neural networks. \textit{J. Comput. Phys.}, 378, 686--707.
\bibitem{karniadakis2021} Karniadakis, G.\,E. et al.\ (2021). Physics-informed machine learning.
\textit{Nature Reviews Physics}, 3, 422--440.
\bibitem{chen2021} Chen, X. et al.\ (2021). Solving inverse stochastic problems via
Fokker-Planck and PINNs. \textit{SIAM J. Sci. Comput.}
\bibitem{krishnapriyan2021} Krishnapriyan, A.\,S. et al.\ (2021). Failure modes in PINNs.
\textit{NeurIPS}, 34.
\end{thebibliography}

\end{document}
